%package list
\documentclass{article}
\usepackage[top=3cm, bottom=3cm, outer=3cm, inner=3cm]{geometry}
\usepackage{graphicx}
\usepackage{url}
\usepackage{hyperref}
\usepackage{array}
\newcolumntype{x}[1]{>{\centering\arraybackslash\hspace{0pt}}p{#1}}
\usepackage{natbib}
\usepackage{pdfpages}
\usepackage{multirow}
\usepackage[T1]{fontenc}
\usepackage{imakeidx}
\usepackage{wrapfig}
\usepackage{caption}
\usepackage{float}

\usepackage{hyperref}
\hypersetup{
    colorlinks=true,
    linkcolor=black,
    filecolor=magenta,
    urlcolor=blue,
    pdftitle={Pruebas Mockito 2},
    pdfpagemode=FullScreen,
}

\urlstyle{same}

\usepackage[normalem]{ulem}
\useunder{\uline}{\ul}{}

\usepackage{minted}

\usepackage{listings}
\usepackage{color, colortbl}
\definecolor{dkgreen}{rgb}{0,0.6,0}
\definecolor{gray}{rgb}{0.5,0.5,0.5}
\definecolor{mauve}{rgb}{0.58,0,0.82}
\definecolor{codebackground}{rgb}{0.95, 0.95, 0.92}
\definecolor{tablebackground}{rgb}{0.0, 0.45, 0.63}

\lstset{
    frame=tb,
    language=bash,
    aboveskip=3mm,
    belowskip=3mm,
    showstringspaces=false,
    columns=flexible,
    basicstyle={\small\ttfamily},
    numbers=none,
    numberstyle=\tiny\color{gray},
    keywordstyle=\color{blue},
    commentstyle=\color{dkgreen},
    stringstyle=\color{mauve},
    breaklines=true,
    breakatwhitespace=true,
    tabsize=3,
    backgroundcolor=\color{codebackground},
}

%%%%%%%%%%%%%%%%%%%%%%%%%%%%%%%%%%%%%%%%%%%%%%%%%%%%%%%%%%%%%%%%%%%%%%%%%%%%
\newcommand{\csemail}{pramirezsa@ulasalle.edu.pe}
\newcommand{\csdocente}{MSc. Luis Pati\~no Herrera}
\newcommand{\cscurso}{Pruebas de Software}
\newcommand{\csuniversidad}{Universidad La Salle}
\newcommand{\csescuela}{Escuela Profesional de Ingenier\'ia de Software}
\newcommand{\cspracnr}{02}
\newcommand{\cstema}{Pruebas con Mockito}
%%%%%%%%%%%%%%%%%%%%%%%%%%%%%%%%%%%%%%%%%%%%%%%%%%%%%%%%%%%%%%%%%%%%%%%%%%%%

\usepackage[english,spanish]{babel}
\usepackage[utf8]{inputenc}
\AtBeginDocument{\selectlanguage{spanish}}
\renewcommand{\figurename}{Figura}
\renewcommand{\refname}{Referencias}
\renewcommand{\tablename}{Tabla}

\usepackage{fancyhdr}
\pagestyle{fancy}
\fancyhf{}
\setlength{\headheight}{30pt}
\renewcommand{\headrulewidth}{1pt}
\renewcommand{\footrulewidth}{1pt}

\fancyhead[L]{\raisebox{-0.2\height}{\includegraphics[width=3cm]{logo_ulasalle (1).png}}}
\fancyhead[R]{\fontsize{7}{7}\selectfont \csuniversidad \\ \csescuela \\ \textbf{\cscurso}}

\fancyfoot[C]{Pruebas de Software}
\fancyfoot[R]{P\'agina \thepage}

\begin{document}

\vspace*{10px}

\begin{center}
    \fontsize{17}{17} \textbf{ Pr\'actica \cspracnr}
\end{center}

\renewcommand{\arraystretch}{1.5}

\begin{table}[h]
\begin{tabular}{|x{4.7cm}|x{4.8cm}|x{4.8cm}|}
\hline
\textbf{DOCENTE} & \textbf{CARRERA} & \textbf{CURSO} \\
\hline
\csdocente & \csescuela & \cscurso \\
\hline
\end{tabular}
\end{table}

\begin{table}[h]
\begin{tabular}{|x{4.7cm}|x{4.8cm}|x{4.8cm}|}
\hline
\textbf{GRUPO} & \textbf{TEMA} & \textbf{DURACI\'ON} \\
\hline
- & \cstema & - \\
\hline
\end{tabular}
\end{table}

\renewcommand{\arraystretch}{1}

\section*{Integrantes}
\begin{itemize}
    \item Patrick Andres Ramirez Santos
\end{itemize}

\section*{Repositorio}
C\'odigo fuente disponible en: \url{https://github.com/patrickram99/pruebas-de-software/tree/main/practica2}

\tableofcontents
\newpage

%% ============================================================
\section{Objetivo}
Implementar pruebas unitarias con JUnit 5 y Mockito para tres estructuras de datos: \texttt{BinaryTree}, \texttt{CircleLinkedList} y \texttt{DynamicArray}. Corregir errores encontrados en el c\'odigo fuente y validar el comportamiento mediante spies y verificaciones.

%% ============================================================
\section{Errores corregidos}

Se identificaron y corrigieron los siguientes errores en el c\'odigo fuente:

\begin{enumerate}
    \item \textbf{BinaryTree.bfs}: El recorrido BFS agregaba el hijo derecho antes que el izquierdo, invirtiendo el orden esperado.
    \item \textbf{CircleLinkedList (constructor)}: El nodo dummy no era circular; \texttt{head.next} era \texttt{null} en vez de apuntar a s\'i mismo.
    \item \textbf{CircleLinkedList.remove}: La validaci\'on de l\'imites usaba \texttt{pos > size} en lugar de \texttt{pos >= size}, permitiendo acceso fuera de rango.
    \item \textbf{DynamicArray.iterator}: Retornaba \texttt{Iterator} sin tipo gen\'erico en vez de \texttt{Iterator<E>}.
    \item \textbf{DynamicArray.next}: La verificaci\'on de l\'imites usaba \texttt{>} en lugar de \texttt{>=}, permitiendo un acceso inv\'alido.
\end{enumerate}

%% ============================================================
\section{Compilaci\'on}

\begin{minted}{bash}
$ mvn compile
\end{minted}

\begin{figure}[H]
\centering
\includegraphics[width=0.9\linewidth]{pruebas/compile.png}
\caption{Compilaci\'on exitosa del proyecto}
\end{figure}

\newpage
%% ============================================================
\section{Ejecuci\'on de pruebas}

\subsection{Todas las pruebas}

\begin{minted}{bash}
$ mvn test
\end{minted}

\begin{figure}[H]
\centering
\includegraphics[width=0.9\linewidth]{pruebas/all_tests.png}
\caption{Ejecuci\'on de las 60 pruebas (BUILD SUCCESS)}
\end{figure}

\subsection{BinaryTreeTest}

\begin{minted}{bash}
$ mvn test -Dtest=BinaryTreeTest
\end{minted}

\begin{figure}[H]
\centering
\includegraphics[width=0.90\linewidth]{pruebas/binary_tree_test.png}
\caption{20 pruebas: find, put, remove y spy con verify}
\end{figure}

\newpage
\subsection{CircleLinkedListTest}

\begin{minted}{bash}
$ mvn test -Dtest=CircleLinkedListTest
\end{minted}

\begin{figure}[H]
\centering
\includegraphics[width=0.90\linewidth]{pruebas/circle_linked_list_test.png}
\caption{17 pruebas: append, remove, toString y spy con verify}
\end{figure}

\subsection{DynamicArrayTest}

\begin{minted}{bash}
$ mvn test -Dtest=DynamicArrayTest
\end{minted}

\begin{figure}[H]
\centering
\includegraphics[width=0.90\linewidth]{pruebas/dynamic_array_test.png}
\caption{23 pruebas: add, remove, iterator, stream y spy con verify}
\end{figure}

%% ============================================================
\newpage
\section{Conclusiones}
\begin{itemize}
    \item Se identificaron y corrigieron 5 errores en las estructuras de datos.
    \item Se implementaron 60 pruebas unitarias con JUnit 5.
    \item Se utiliz\'o Mockito (spy, verify, clearInvocations) para validar interacciones internas.
    \item Todas las pruebas pasan exitosamente tras las correcciones.
\end{itemize}

\end{document}
